
\documentclass[12pt]{article}

\usepackage{geometry}
\usepackage{titling}
\usepackage{amsmath,amssymb}
\usepackage{fontspec}
\usepackage{unicode-math}
\usepackage{array}
\usepackage{colortbl}
\usepackage[dvipsnames]{xcolor}
\usepackage[most]{tcolorbox}

\geometry{a4paper, top=2.5cm, bottom=2.5cm, left=2.6cm, right=2.5cm}
\setlength{\droptitle}{-3cm}

\setmonofont{DejaVu Sans Mono}
\setmainfont{Libertinus Serif}
\setsansfont{Libertinus Sans}
\setmathfont{STIXTwoMath-Regular.otf}

\newcommand{\MVec}[1]{\mathbf{#1}}

\title{Application of Euclidean Distance \\ in \\ K-Nearest Neighbor Regression}
\author{Devi Prasad M}
\date{August 2026}

\begin{document}
\pagecolor{yellow!30!brown!40}

\maketitle

\section*{Euclidean Distance -- An Application: KNN Regression}

\subsection*{A 4D Example: Predicting House Prices}
Consider a 4-dimensional feature space where we want to predict the sale price of a house based on its characteristics. Our features are:
\begin{itemize}
    \item $x^{(1)}$: Living Area (in $1,000$ sq ft) -- \textit{Continuous}
    \item $x^{(2)}$: Distance to City Center (kms) -- \textit{Continuous}
    \item $x^{(3)}$: Age of Property (years) -- \textit{Integer}
    \item $x^{(4)}$: Number of Bedrooms -- \textit{Integer}
\end{itemize}
Our target $y$ is the \textbf{Sale Price (in INR Crores)}.

\vspace{0.5cm}
\noindent Suppose a new house comes on the market. Our query point is a $2,000$ sq ft home, $1.0$ km from the center, $10$ years old, with $3$ bedrooms:
\[
\MVec{x}_q = (2.0, 1.0, 10, 3)
\]

\noindent Suppose our historical training dataset contains the following $12$ recent sales:
\[
\begin{array}{c|cccc|c}
\text{Property} & x^{(1)} \text{ (Area)} & x^{(2)} \text{ (Dist)} & x^{(3)} \text{ (Age)} & x^{(4)} \text{ (Beds)} & y_i \text{ (Price)} \\ \hline
\MVec{x}_1 & 1.5 & 1.5 & 12 & 3 & 3.5 \\
\MVec{x}_2 & 2.2 & 0.5 & 8  & 4 & 5.2 \\
\MVec{x}_3 & 2.0 & 1.0 & 9  & 3 & 4.8 \\
\MVec{x}_4 & 1.8 & 1.2 & 10 & 3 & 4.6 \\
\MVec{x}_5 & 2.5 & 2.0 & 5  & 4 & 6.0 \\
\MVec{x}_6 & 1.2 & 3.0 & 20 & 2 & 2.5 \\
\MVec{x}_7 & 2.1 & 0.8 & 11 & 3 & 4.9 \\
\MVec{x}_8 & 3.0 & 5.0 & 2  & 5 & 7.0 \\
\MVec{x}_9 & 1.9 & 1.5 & 10 & 2 & 4.1 \\
\MVec{x}_{10} & 2.0 & 0.5 & 15 & 3 & 4.3 \\
\MVec{x}_{11} & 2.3 & 1.1 & 10 & 4 & 5.4 \\
\MVec{x}_{12} & 1.7 & 1.0 & 8  & 3 & 4.4
\end{array}
\]
We choose $K=3$.

\subsection*{Step 1: Measure the Distance to Every Point}
The 4-dimensional Euclidean distance to any training point $\MVec{x}_i$ is:
\[
\mathbf{dist}(\MVec{x}_q,\MVec{x}_i)
=
\sqrt{
(2.0-x_i^{(1)})^2+
(1.0-x_i^{(2)})^2+
(10-x_i^{(3)})^2+
(3-x_i^{(4)})^2
}
\]
We calculate this geometric distance to \textit{every} one of the $12$ properties:
\[
\begin{array}{c|c|c@{\qquad}c|c|c}
\text{Property} & y_i \text{ (Price)} & \text{Distance}
& \text{Property} & y_i \text{ (Price)} & \text{Distance} \\ \hline
\MVec{x}_1 & 3.5 & 2.12 & \MVec{x}_7 & 4.9 & 1.02 \\
\MVec{x}_2 & 5.2 & 2.30 & \MVec{x}_8 & 7.0 & 9.22 \\
\MVec{x}_3 & 4.8 & 1.00 & \MVec{x}_9 & 4.1 & 1.12 \\
\MVec{x}_4 & 4.6 & 0.28 & \MVec{x}_{10} & 4.3 & 5.02 \\
\MVec{x}_5 & 6.0 & 5.22 & \MVec{x}_{11} & 5.4 & 1.05 \\
\MVec{x}_6 & 2.5 & 10.28 & \MVec{x}_{12} & 4.4 & 2.02
\end{array}
\]
\[
\boxed{
\text{Measure: Every training point is compared with the query.}
}
\]

\subsection*{Step 2: Select the $K$ Nearest Neighbors}
Since $K=3$, we select the three properties with the smallest Euclidean distances:
\[
\begin{array}{c|c|c}
\text{Neighbor} & \text{Target } y_i & \mathbf{dist}(\MVec{x}_q,\MVec{x}_i) \\ \hline
\MVec{x}_4=(1.8, 1.2, 10, 3) & 4.6 & 0.28 \\
\MVec{x}_3=(2.0, 1.0, 9, 3) & 4.8 & 1.00 \\
\MVec{x}_7=(2.1, 0.8, 11, 3) & 4.9 & 1.02
\end{array}
\]
Therefore, our nearest neighborhood is:
\[
N_3(\MVec{x}_q)
=
\left\{
(\MVec{x}_4, 4.6),
(\MVec{x}_3, 4.8),
(\MVec{x}_7, 4.9)
\right\}.
\]

\subsection*{Step 3: Predict}
The KNN regression prediction is the arithmetic mean of the target prices of these three most similar houses:
\[
\hat{y}(\MVec{x}_q)
=
\frac{4.6 + 4.8 + 4.9}{3}
=
\frac{14.3}{3}
\approx
4.77.
\]
The KNN estimate for the unknown market value of our query house is:
\[
\boxed{
\text{Estimated Price } \approx \text{INR } 4.77 \text{ Crores}
}
\]

\newpage

\subsection*{A Critical Flaw: The Need for Feature Scaling}
While the example above perfectly demonstrates the \textit{mechanics} of KNN, there is a hidden mathematical flaw in applying Euclidean distance directly to raw, unscaled real-world data.

Consider the differences in scale among our features:
\begin{itemize}
    \item A difference of $5$ years in \textbf{Age} contributes $(5)^2 = 25$ to the distance sum.
    \item A difference of $2$ kms in \textbf{Distance} contributes only $(2)^2 = 4$ to the distance sum.
\end{itemize}

Even if distance to the city center has a more massive impact on actual property value, the \textbf{Age} feature mathematically dominates the Euclidean distance simply because its numerical values are larger. In effect, we are comparing apples to oranges.

\vspace{0.3cm}
\noindent\textbf{The Solution: Standardization} \\
To prevent variables with larger scales from biasing the distance metric, distance-based algorithms require data preprocessing. We must normalize or standardize the features so they contribute equally to the geometry.

A standard approach is \textit{Z-score standardization}, which transforms every feature to have a mean of $0$ and a standard deviation of $1$:
\[
x^{(j)}_{\text{scaled}} = \frac{x^{(j)} - \mu_j}{\sigma_j}
\]
where $\mu_j$ is the mean and $\sigma_j$ is the standard deviation of feature $j$.

\end{document}